\documentclass[11pt,a4paper]{article}
\newcommand{\Rodzaj}{Szablon sprawozdania}
\usepackage{fontspec}
\setmainfont{texgyretermes-regular.otf}[BoldFont=texgyretermes-bold.otf,ItalicFont=texgyretermes-italic.otf]
\usepackage[polish]{babel}
\usepackage{amsmath,amssymb,geometry,array,tabularx,graphicx,caption,fancyhdr}
\usepackage[table]{xcolor}
\geometry{left=18mm,right=18mm,top=19mm,bottom=19mm,headheight=14pt,headsep=6mm}
\setlength{\parindent}{0pt}\setlength{\parskip}{5pt}
\setlength{\emergencystretch}{2em}
\renewcommand{\arraystretch}{1.3}\setlength{\tabcolsep}{5pt}
\setlength{\fboxsep}{6pt}
\setlength{\abovedisplayskip}{8pt}\setlength{\belowdisplayskip}{8pt}
\setlength{\abovedisplayshortskip}{5pt}\setlength{\belowdisplayshortskip}{5pt}
\widowpenalty=10000\clubpenalty=10000
\raggedbottom
\newcolumntype{Y}{>{\raggedright\arraybackslash}X}
\newcolumntype{C}{>{\centering\arraybackslash}X}
\newcommand{\naglowek}{\rowcolor{gray!12}}
\newcommand{\odstep}{\par\smallskip}
\captionsetup{font=small,labelfont=bf,justification=justified,singlelinecheck=false}
\pagestyle{fancy}\fancyhf{}\fancyhead[L]{\small Ćwiczenie 11 — moduł Younga}
\fancyhead[R]{\small\Rodzaj}\fancyfoot[C]{\small\thepage}
\renewcommand{\headrulewidth}{0.3pt}
\makeatletter
\renewcommand\section{\@startsection{section}{1}{0pt}{10pt}{4pt}{\large\bfseries}}
\renewcommand\subsection{\@startsection{subsection}{2}{0pt}{8pt}{4pt}{\normalsize\bfseries}}
\makeatother
\newcommand{\pole}{\rule{23mm}{0.3pt}}
\newcommand{\miejsce}[2]{\par\medskip\noindent\fbox{\begin{minipage}[t][#1][t]{\dimexpr\linewidth-2\fboxsep-2\fboxrule\relax}\small #2\end{minipage}}\par\medskip}
\newcommand{\metryka}{%
\begingroup\small\setlength{\parskip}{0pt}\renewcommand{\arraystretch}{1.1}
\begin{tabularx}{\linewidth}{|p{29mm}|Y|p{18mm}|p{18mm}|p{19mm}|}\hline
\textbf{Wydział}\newline Informatyki & \textbf{Imię i nazwisko}\newline 1. Marcin Podyma\newline 2. Paweł Mierzwa & \textbf{Rok}\newline 2026 & \textbf{Grupa}\newline 3 & \textbf{Zespół}\newline 1\\\hline
\end{tabularx}\par\nointerlineskip
\begin{tabularx}{\linewidth}{|p{29mm}|Y|p{19mm}|}
\parbox[t]{29mm}{\centering\textbf{PRACOWNIA\\FIZYCZNA\\WFiIS AGH}} & \textbf{Temat:}\newline Wyznaczanie modułu Younga metodą statyczną & \textbf{Nr}\newline\textbf{ćwiczenia}\newline 11\\\hline
\end{tabularx}\par\nointerlineskip
\begin{tabularx}{\linewidth}{|p{29mm}|Y|Y|Y|Y|p{19mm}|}
\textbf{Data wykonania}\newline\strut & \textbf{Data oddania}\newline\strut & \textbf{Zwrot do popr.} & \textbf{Data oddania} & \textbf{Data zaliczenia} & \textbf{OCENA}\\\hline
\end{tabularx}\endgroup\medskip}

\newcommand{\przyrzadopis}[1]{%
\par\medskip\noindent
\begin{minipage}[c]{0.44\linewidth}
\centering\includegraphics[width=\linewidth,height=104mm,keepaspectratio]{#1}
\end{minipage}\hfill
\begin{minipage}[c]{0.51\linewidth}
\raggedright
\captionof{figure}{Przyrząd do statycznego pomiaru modułu Younga.}
\medskip
\textbf{A — górny uchwyt}\par
Mocuje górny koniec badanego drutu.
\par\smallskip\textbf{B — dolne mocowanie}\par
Łączy drut z układem obciążenia.
\par\smallskip\textbf{C — dźwignia}\par
Przenosi przemieszczenie drutu na czujnik.
\par\smallskip\textbf{D — czujnik mikrometryczny}\par
Wskazuje przemieszczenie dwukrotnie większe od wydłużenia drutu.
\par\smallskip Szalka z odważnikami obciąża dolne mocowanie.
\medskip\par{\footnotesize Rysunek: instrukcja \texttt{11\_wykon.pdf}, s.~2.}
\end{minipage}\par\medskip}
\newcommand{\teoria}{%
W zakresie liniowej sprężystości prawo Hooke'a wiąże naprężenie
$\sigma$ z odkształceniem względnym $\varepsilon$:
\[
\boxed{\sigma=E\varepsilon,\quad\sigma=\frac{F}{S},\quad
\varepsilon=\frac{\Delta l}{l},\quad S=\frac{\pi d^2}{4}.}
\]
Moduł Younga $E$ ma jednostkę Pa. Zależność wydłużenia od siły
w przyjętym zakresie opisujemy prostą ze swobodnym wyrazem wolnym:
\[
\boxed{\Delta l=aF+b,\qquad E=\frac{4l}{\pi d^2a},\qquad F=mg.}
\]
Wzór na $E$ wymaga jednostek SI: $l,d$ w m oraz $a$ w m/N.
Dla dopasowania względem masy $\Delta l=cm+b$ zamiast siły
należy użyć $E=4lg/(\pi d^2c)$; nie wolno zamieniać współczynników $a$ i $c$.
}
\newcommand{\dzwignia}{%
Czujnik jest sprzężony z drutem przez dźwignię. Odległość drutu
od osi obrotu jest dwa razy mniejsza niż odległość punktu styku czujnika,
zatem jego wskazanie jest dwa razy większe od wydłużenia drutu.
Po prawidłowym wyzerowaniu:
\[
\boxed{\Delta l_{\uparrow}=\frac{z_{\uparrow}}2,\quad
\Delta l_{\downarrow}=\frac{z_{\downarrow}}2,\quad
\overline{\Delta l}=\frac{z_{\uparrow}+z_{\downarrow}}4.}
\]
Jeśli zero nie jest równe 0, przed użyciem wzoru odejmij odpowiedni
odczyt zerowy. Nie dziel wskazań przez 2 ponownie w regresji.
}
\newcommand{\tabelapomiarow}{%
\begingroup\small\renewcommand{\arraystretch}{1}
\begin{tabularx}{\linewidth}{|c|C|C|C|C|Y|}\hline
\naglowek\textbf{Nr} & \textbf{Masa całkowita} [kg] & \textbf{Czujnik $z_\uparrow$} [mm] &
\textbf{Czujnik $z_\downarrow$} [mm] & \textbf{Zero / uwagi} & \textbf{Uwagi do punktu}\\\hline
\csname @@input\endcsname{tabela_pusta.tex}
\end{tabularx}\endgroup}
\newcommand{\tabelaopracowania}{%
\begingroup\small\renewcommand{\arraystretch}{1}
\begin{tabularx}{\linewidth}{|c|C|C|C|C|C|}\hline
\naglowek\textbf{Nr} & $m$ [kg] & $F=mg$ [N] &
$z_\uparrow$ [mm] & $z_\downarrow$ [mm] & $\overline{\Delta l}$ [mm]\\\hline
\csname @@input\endcsname{tabela_pusta.tex}
\end{tabularx}\endgroup}
\newcommand{\wymiary}[1]{%
\textbf{#1}\quad Materiał: \pole\quad Długość $l=$ \pole\ m.
\par\smallskip
\begin{tabularx}{\linewidth}{|Y|C|C|C|}\hline
\naglowek\textbf{Średnica w trzech miejscach} & $d_1$ [mm] & $d_2$ [mm] & $d_3$ [mm]\\\hline
Pomiar przy obciążeniu 2 kg & & &\rule{0pt}{9mm}\\\hline
\end{tabularx}\par\smallskip}

\begin{document}\metryka
\section{Cel ćwiczenia}
Wyznaczenie modułu Younga dwóch metali metodą statyczną na podstawie
zależności wydłużenia drutu od obciążenia.
\section{Podstawy teoretyczne}\teoria
\section{Aparatura i przebieg pomiarów}\dzwignia
Układ obejmuje statyw z uchwytami drutu i szalką, dźwignię,
czujnik mikrometryczny, odważniki, mikrometr i przymiar długości.
\clearpage
\subsection{Przyrząd pomiarowy}
\przyrzadopis{Rysunki/przyrzad.pdf}
Przymiarem wyznacz długość badanego odcinka drutu, a mikrometrem
średnicę w trzech miejscach. W opisie podaj materiały, zakresy obciążeń,
krok masy, zerowanie oraz sposób odczytu przy dokładaniu i zdejmowaniu odważników.
\miejsce{25mm}{Rzeczywisty przebieg: materiały, zakresy mas, krok, zerowanie,
pomiar długości i średnicy, obciążanie i odciążanie; odstępstwa od instrukcji.}

\clearpage\section{Wyniki pomiarów}
\subsection{Drut A}\wymiary{Parametry próbki}
\textbf{Niepewności typu B:} $u(l)=$ \pole\ m,
$u(d)=$ \pole\ mm. Podaj sposób ich wyznaczenia.

W tabeli zachowaj odczyty obu kierunków. Siłę $F=mg$ i średnie
wydłużenie $(z_\uparrow+z_\downarrow)/4$ oblicz po przepisaniu danych.
\par\medskip\tabelaopracowania
\clearpage\subsection{Drut B}\wymiary{Parametry próbki}
\textbf{Niepewności typu B:} $u(l)=$ \pole\ m,
$u(d)=$ \pole\ mm. Podaj sposób ich wyznaczenia.

Zapisz zerowanie i oznacz odczyt maksymalny, jeśli wykorzystano go
dwukrotnie. Nie dopisuj punktu $(0,0)$ jako pomiaru, gdy go nie wykonano.
\par\medskip\tabelaopracowania

\clearpage\section{Regresja i wyznaczenie modułu}
Oblicz $\bar d=(d_1+d_2+d_3)/3$. Zgodnie z instrukcją niepewność
średnicy oceń metodą typu B, nie z rozrzutu trzech średnic.
Dla działki $q$ i założenia błędu zaokrąglenia w $[-q/2,q/2]$
otrzymujemy $u_B=q/\sqrt{12}$. Jeśli podano granicę błędu
$\pm\Delta$ z rozkładem prostokątnym, użyj $u_B=\Delta/\sqrt3$.
Nie utożsamiaj działki z granicą błędu bez wyjaśnienia.

Pokaż wszystkie zmierzone punkty. Punkty pominięte w dopasowaniu
zaznacz i uzasadnij wybór zakresu, bez dopasowywania go do wartości tablicowej.
Dopasuj $\Delta l=aF+b$; podaj $a$, $u(a)$, $b$ i zakres regresji.
\miejsce{65mm}{Wykres drutu A: $\Delta l(F)$, obie osie z jednostkami,
punkty użyte i pominięte, prosta, zakres i parametry dopasowania.}
\miejsce{65mm}{Wykres drutu B: $\Delta l(F)$, obie osie z jednostkami,
punkty użyte i pominięte, prosta, zakres i parametry dopasowania.}
\[
\boxed{E=\frac{4l}{\pi\bar d^2a},\qquad
a\,[\mathrm{m/N}]=10^{-3}a\,[\mathrm{mm/N}].}
\]

\clearpage\section{Niepewności i porównanie z odniesieniem}
Przy niezależnych $l,d,a$ i pomijalnej niepewności siły, jak w instrukcji:
\[
\boxed{\frac{u_c(E)}E=\sqrt{
\left(\frac{u(l)}l\right)^2+
\left(2\frac{u(d)}d\right)^2+
\left(\frac{u(a)}a\right)^2},\qquad U(E)=k\,u_c(E).}
\]
Współczynnik 2 przy $u(d)/d$ wynika z potęgi $d^2$.
Podaj $k$; dla $k=2$ nie deklaruj automatycznie dokładnie 95\% pokrycia.
Jeśli używasz wykresu względem masy i niepewnego $g$, dodaj także wkład $u(g)/g$.
\begin{tabularx}{\linewidth}{|Y|C|C|C|C|}\hline
\naglowek\textbf{Materiał} & $E$ [GPa] & $u_c(E)$ [GPa] & $U(E)$ [GPa] & \textbf{Odniesienie} [GPa]\\\hline
Drut A & & & &\rule{0pt}{15mm}\\\hline
Drut B & & & &\rule{0pt}{15mm}\\\hline
\end{tabularx}
\miejsce{27mm}{Porównanie: źródło odniesienia, przedział $E\pm U$,
zgodność lub jej brak. Dyskusja wpływu średnicy, zerowania, wyrazu wolnego,
histerezy i wyboru zakresu. Nie przypisuj przyczyny bez dowodów.}
\section{Wnioski}
\begin{tabularx}{\linewidth}{|p{42mm}|Y|}\hline
\naglowek\textbf{Zagadnienie} & \textbf{Wynik i interpretacja}\\\hline
Wyniki obu materiałów &\rule{0pt}{14mm}\\\hline
Porównanie tablicowe &\rule{0pt}{14mm}\\\hline
Zakres liniowy i powrót do zera &\rule{0pt}{14mm}\\\hline
Dominujące niepewności / ograniczenia &\rule{0pt}{14mm}\\\hline
\end{tabularx}
\end{document}
